\section{引言}
\label{sec:intro}

维护脆弱的 UI 测试套件是现代 Web 工程中持续存在的维护成本。一次 CSS 类名重命名、一次 React 组件的重新嵌套，或是一次行为不变的处理函数重构，都足以让数十个端到端（E2E）Playwright 测试断裂。无论是工业界工具（Healenium、mabl）还是日益增长的学术工作~\cite{joseph2026zerocost, mouraReproBreak2026, healenium}，给出的自然回应都是\emph{自愈}: 在定位器失效时，自动将其重写为指向新 DOM 中等价元素的形式。

这一承诺是真实的，最近的文献也使其看起来已经接近被解决。Joseph~\cite{joseph2026zerocost} 通过一个确定性可访问性阶梯在单个 demo 上报告了 100\% 通过率；当代的 LLM 定位器修复器也用各种一致性守门规则报告了较高的修复成功率。但这些数字共享一个结构性的盲区: 它们衡量的是\emph{打过补丁的测试是否通过}，而不是\emph{补丁是否指向开发者真正想要的元素}。一个总能找到\emph{某个}可解析选择器（任何属性近似的元素都行）的修复器是沉默地危险的，因为每一次"通过但错指"的修复都会遮蔽一次本应被作为回归暴露出来的真实 UI 变更。

我们在 ReproBreak~\cite{mouraReproBreak2026} 数据集中的真实 Playwright 失效用例上直接测量这种失败模式。在 koenig 的 75 个可达失效用例上，结果令人不安:

\begin{itemize}
\item \textbf{F1 —— 隐性误修复是默认行为。}当 ground truth 目标元素被强制从候选集中移除时，一个朴素的 LLM 修复器 (\texttt{qwen2.5-coder:7b}~\cite{qwenCoderTech2024}) 在 \textbf{78.7\%} 的用例中给出了语法正确但元素错误的选择器（59/75）；仅 17.3\% 的时候它能正确放弃。
\item \textbf{F2 —— 软提示放弃指令无效。}向系统提示中追加一条明确的"如无候选契合则放弃"守门规则后，结果数字\emph{完全一致}（两个版本都是 59/75，$\Delta=0$）。模型把这条规则当作装饰物。
\item \textbf{F3 —— 跨应用泛化被锚点文化的间接性所限制。}另一个 React+Playwright 代码库（\texttt{payloadcms/payload}，319 个失效配对）使用 BEM 风格的模板字面量类名，而当前的静态抽取器无法对其求值，原始可达率从 koenig 的 80.6\% 降到 3.8\%。我们将其诊断为一个明确的、可修复的盲区，而非对方法的否证。
\end{itemize}

\noindent 前两个发现给出了\emph{为无人值守自愈流水线的下一层评估行为重放验证}的直接经验性动机；我们尚未在重放下测量端到端的遮蔽率（\S\ref{sec:roadmap}，R1+R3）。第三个发现给出了一个可证伪的跨应用泛化论断: 任何在单一应用上 benchmark 的定位器修复工具都很可能高估其泛化能力，因为锚点文化（\texttt{data-testid} vs \texttt{data-kg-*} vs BEM 模板字面量）在不同 React 代码库间存在显著差异。

此外，我们介绍 HealReact 的 L1 基底 —— 一个 AST 抽取器加一个带确定性后置校准的 LLM 派生意图标注器 —— 它静态可达 koenig 93 个失效目标中的 80.6\%；在其之上，一个小型本地 L3 修复器在可达用例上正确回答了 77.3\%（在完整数据集上为 62.4\% 的静态修复代理\footnote{我们全文使用"静态修复代理"而非"端到端修复": 解析器在元素层面的精确一致性与 Playwright 通过率\emph{相关}，但既不是其下界也不是其上界（静态精确匹配的选择器可能在 visibility/timing 约束下运行时失败；静态不匹配的选择器也可能在重渲染 DOM 中运行时通过）。两者之间的换算需要基于 Docker 的重放，我们将其列为优先 roadmap（\S\ref{sec:roadmap}）。}）。两个廉价确定性杠杆 —— testId 加权检索与强锚点后置过滤 —— 带来 5 个新正确用例、零回归的帕累托改进，表明确定性后处理比提示工程更可靠地起到安全网作用。

\paragraph{贡献。}
（i）在真实 Playwright 数据集上对隐性误修复的诊断性测量，以及对软提示放弃指令无效的展示（\S\ref{sec:findings}）；
（ii）HealReact 的 L1 基底作为一个可证伪、可本地运行的静态层（\S\ref{sec:l1}）；
（iii）两个对朴素 LLM 修复器形成帕累托改进的确定性安全杠杆（\S\ref{sec:findings}）。

78.7\% 是上界（对抗性构造）；在 koenig 真正不可达的 18/93 用例上换算得到无人值守工作负载下约 15\% 的隐性误修复率。我们（\S\ref{sec:discussion}）认为这已经高到足以拒绝在 CI 中无人值守地自动提交 LLM 提出的选择器补丁。
